\section{Introduction}
\label{sec:intro}

Multi-agent reinforcement learning systems deployed in the real world rarely operate in lockstep. A warehouse robot that plans routes every second coordinates with a pick-and-place arm that decides every 100 milliseconds; a fleet manager issuing high-level directives operates on a fundamentally different timescale than the vehicles executing them. This temporal heterogeneity is not a nuisance to be engineered away---it is an inherent property of complex multi-agent systems. Yet the dominant paradigm in cooperative \marl assumes synchronous agents with identical decision frequencies \citep{rashid2018qmix, yu2022mappo, foerster2018coma}.

A tempting hypothesis is that temporal heterogeneity itself creates pressure for hierarchical organization. Agents with longer decision horizons might naturally learn to serve as implicit coordinators, developing value functions that model and anticipate the behavior of faster agents. If true, this would provide theoretical grounding for designing scalable \marl systems without hand-crafted coordination hierarchies---temporal diversity alone would suffice. This intuition draws support from hierarchical RL, where architectures like \feudal Networks \citep{vezhnevets2017feudal} and the Option-Critic framework \citep{bacon2017option} achieve coordination precisely through multi-timescale decision-making.

However, these hierarchical RL methods \emph{impose} structure by design. Whether analogous structure \emph{emerges} from temporal heterogeneity in multi-agent settings remains an open question. Recent work on credit assignment in cooperative \marl has made progress on the related but distinct problem of decomposing joint rewards across both agents and timesteps. \tartwo \citep{kapoor2024tar2} provides a principled framework for joint agent-temporal credit decomposition, and \stas \citep{chen2024stas} uses dual transformers for spatial-temporal return decomposition. Yet both assume synchronous agents with uniform horizons, leaving the heterogeneous case unexplored.

We address this gap directly. We train Independent PPO (\ippo) agents---intentionally the simplest multi-agent algorithm---with heterogeneous action frequencies across three cooperative gridworld environments that vary in their synchronization demands. We use \ippo because any coordination must emerge from the shared reward signal alone, without explicit inter-agent information sharing that could confound the analysis. We quantify emergent hierarchy through mutual information between agent value functions at multiple time lags, looking for the directional asymmetry (slow $\to$ fast $>$ fast $\to$ slow) that would indicate hierarchical credit assignment.

Our main finding is negative: \textbf{temporal heterogeneity does not induce emergent hierarchical credit assignment}. Across all three environments and 45 experimental runs, we observe no directional mutual information asymmetry between slow and fast agents' value functions. Instead, heterogeneity creates task-dependent performance effects: a 45\% advantage in complementary-role tasks ($p < 0.0001$) but a 17\% disadvantage in synchronization-dependent tasks ($p = 0.001$). These results suggest that the shared reward signal in cooperative \marl is insufficient to drive hierarchical organization without explicit architectural inductive biases.

Our contributions are as follows:
\begin{itemize}[leftmargin=*,itemsep=0pt,topsep=0pt]
    \item We design three cooperative gridworld environments with configurable agent action frequencies to systematically study the effect of temporal heterogeneity on multi-agent coordination.
    \item We conduct a controlled study across 45 runs with mutual information analysis, demonstrating that hierarchical credit assignment does not spontaneously emerge from heterogeneous decision horizons in \ippo.
    \item We identify a task-structure dependence: temporal heterogeneity helps when roles are naturally complementary but hurts when synchronization is required, clarifying when heterogeneity is beneficial versus harmful.
    \item We show through ablation that temporal context features significantly improve coordination in synchronization-critical tasks, providing evidence that agents \emph{can} learn to use temporal awareness even though hierarchy does not emerge.
\end{itemize}
